\section{Trajectory-aware machine design and the unified framework}
\label{sec:dynamic-design}

The dynamic geometry can be read backward just like the static value
functions.  For representative transitions \(A_k\rightarrow B_k\), define
\[
 \gls{sym:tdyn}=\min_{\gamma,u}t_f,
 \qquad
 E_{\min,k}(\boldsymbol\theta)=\min_{\gamma,u}E_{\mathrm{loss}},
 \tag{34.1}
\]
under the declared converter, torque, current, and thermal constraints.  The
dynamic design region is
\[
 \Theta_{\mathrm{dyn}}=\bigcap_k
 \{\boldsymbol\theta:t_{\min,k}\leq t_{\mathrm{req},k},
 E_{\min,k}\leq E_{\mathrm{req},k}\}.
 \tag{34.2}
\]
It intersects, rather than replaces, the static torque--speed design region.

Useful screening descriptors include
\[
 \lambda_{\min}(\matr W_c),\quad
 \lambda_{\max}(\matr W_c),\quad
 \kappa(\matr W_c),\quad
 \frac{\tau_f}{\tau_s},\quad
 \lambda_{\min}(\matr G),\quad
 \lambda_{\max}(\matr G),
 \tag{34.3}
\]
all tied to a stated horizon, speed, voltage set, and operating region.  The
Gramian measures finite-horizon control-energy reachability; the local metric
measures instantaneous voltage-to-current speed.  Neither alone certifies a
constrained transition, but both cheaply reject dynamically poor candidates.

\begin{figure}[tb]
\centering
\begin{tikzpicture}[font=\small]
 \begin{axis}[width=.76\linewidth,height=.47\linewidth,axis lines=middle,
  xmin=0,xmax=4.5,ymin=0,ymax=4.0,ticks=none,
  xlabel={static capability $C_M$},ylabel={dynamic reachability $D_M$},clip=false]
  \addplot[name path=static,blue,very thick,domain=1.6:4.2] ({x},{.35});
  \addplot[name path=dyn,BrickRed,very thick,domain=.4:4.2] ({.55},{x});
  \path[name path=top] (axis cs:.55,3.7)--(axis cs:4.2,3.7);
  \path[name path=right] (axis cs:4.2,.35)--(axis cs:4.2,3.7);
  \addplot[green!20] fill between[of=static and top,soft clip={domain=1.6:4.2}];
  \fill[blue!12] (axis cs:1.6,.35) rectangle (axis cs:4.2,3.7);
  \draw[blue,very thick] (axis cs:1.6,.35)--(axis cs:1.6,3.7);
  \draw[BrickRed,very thick] (axis cs:.55,1.35)--(axis cs:4.2,1.35);
  \fill[ForestGreen!70!black] (axis cs:2.9,2.35) circle (2.5pt) node[above right] {robust candidate};
  \fill[Orange] (axis cs:3.1,.8) circle (2.5pt) node[below right,align=left] {statically good,\\dynamically slow};
  \fill[Purple] (axis cs:1.05,2.55) circle (2.5pt) node[above left,align=right] {fast, but insufficient\\steady capability};
  \node[blue] at (2.4,3.55) {$C_M\geq C_{\rm req}$};
  \node[BrickRed,anchor=west] at (axis cs:.65,1.48) {$t_{\min}\leq t_{\rm req}$};
 \end{axis}
\end{tikzpicture}
\caption{Trajectory-aware design intersects static and dynamic capability
regions.  A machine can pass either screening layer and still fail the other.}
\label{fig:dynamic-design-region}
\end{figure}



The design trade-off is physical.  Low rotor resistance improves steady-state
field copper efficiency but, with large field inductance and limited field
voltage, can increase the field time constant.  Large inductance may reduce
ripple yet shrink current acceleration for a fixed DC link.  Strong mutual
coupling improves torque per field ampere while rotating the dynamic axes and
approaching the positive-energy boundary
\(L_dL_f-L_m^2>0\).  A statically efficient machine can therefore be
dynamically sluggish.  Machine-control co-design must price both.

For a physical geometry \(\vect g\), the extended forward map is no longer
only \(\boldsymbol\theta=\mathcal M(\vect g)\).  It also supplies differential
flux maps, thermal resistances, converter ratings, and uncertainty envelopes.
An inexpensive workflow is:
\begin{enumerate}
 \item apply the static capability filters of Part~II;
 \item evaluate local metric and Gramian descriptors at representative points;
 \item solve the reduced transition set (34.1) for survivors;
 \item verify saturation, losses, demagnetization, mechanics, and thermal
 dynamics with FEA-coupled simulation.
\end{enumerate}
The analytic descriptors prioritize expensive simulations; they do not
replace them.

\subsection{Three geometries, one co-design problem}

\begin{figure}[tb]
\centering
\begin{tikzpicture}[font=\small,node distance=19mm and 18mm]
 \node[draw,rounded corners,fill=blue!10,minimum width=35mm,minimum height=16mm,align=center] (static)
 {static geometry\\where to operate?};
 \node[draw,rounded corners,fill=BrickRed!10,minimum width=39mm,minimum height=16mm,align=center,right=of static] (dynamic)
 {dynamic geometry\\how to move?};
 \node[draw,rounded corners,fill=ForestGreen!12,minimum width=39mm,minimum height=16mm,align=center,below right=13mm and -9mm of static] (design)
 {design geometry\\which machine?};
 \draw[<->,very thick] (static)--node[above,font=\scriptsize] {endpoints}(dynamic);
 \draw[<->,very thick] (static)--node[left,font=\scriptsize,align=right] {capability\\value functions}(design);
 \draw[<->,very thick] (dynamic)--node[right,font=\scriptsize,align=left] {transition\\descriptors}(design);
 \node[below=11mm of design,align=center,text width=10cm]
 {geometry and materials $\rightarrow$ flux/loss/dynamics $\rightarrow$ feasible points and paths $\rightarrow$ controller};
\end{tikzpicture}
\caption{The unified machine-control co-design framework.  Static optima are
the endpoints of dynamic paths; both impose inverse constraints on machine
parameters and physical geometry.}
\label{fig:three-geometries}
\end{figure}



The completed framework has three linked questions:
\[
 \boxed{\begin{aligned}
 &\text{Where should the machine operate?}\\
 &\text{How should it move between operating points?}\\
 &\text{Which machine parameters make both possible?}
 \end{aligned}}
 \tag{34.4}
\]
The resulting method hierarchy is evidence-based.  Constant linear dynamics
with ellipsoidal authority admits whitening, exact Gramians, and in special
cases closed geodesics.  Hard converter and path constraints lead to
piecewise-analytic or boundary-following arcs.  PWA saturation suggests
piecewise-flat refraction.  A fully nonlinear machine supports local metrics,
short-horizon direct optimization, or an offline value field; intrinsic
curvature explains why no global coordinate transform need straighten all
paths.  Geometric optimal control provides the common language
\parencite{agrachev2004}, not a mandate to force every case into one metric.

Open questions are now concrete: certified error bounds for PWA refraction,
viscosity-solution regularity of converter-polytope HJB equations, nonsmooth
sensitivities of \(t_{\min,k}\) at switching changes, robust curvature
estimation from flux maps, minimal-memory value-function representations, and
whether a finite set of transition descriptors can predict full drive-cycle
dynamic capability.
